\section{\label{sec:intro}Introduction}

The Standard Model (SM) of particle physics, describing the fundamental interactions of Nature, is among the most successful theories of physics.
However, the SM alone is unable to provide a satisfactory answer to several open questions of both observational and theoretical nature. 
It should therefore be regarded as the low-energy remnant of a more complete ultraviolet theory entailing new dynamics emerging at some large 
--- yet unknown --- energy scale $\Lambda$.

SM extensions including light pseudoscalars, such as the so-called 
Axion-Like Particles (ALPs)~\cite{Jaeckel:2010ni,Marsh:2015xka,Irastorza:2018dyq,DiLuzio:2020wdo}, 
are among the most interesting and studied scenarios.
Their lightness, compared to the scale $\Lambda$, can be easily motivated if 
they are the pseudo-Nambu-Goldstone bosons of some spontaneously broken global symmetry.
ALPs can elegantly address several open questions in particle physics such as the strong CP~\cite{Peccei:1977hh,Peccei:1977ur,Weinberg:1977ma,Wilczek:1977pj} and flavor problems~\cite{Davidson:1981zd,Wilczek:1982rv,Berezhiani:1989fp,Calibbi:2016hwq,Greljo:2024evt}, as well as the stability of the electroweak scale~\cite{Graham:2015cka}. 
Moreover, they can be regarded as being natural dark matter candidates~\cite{Abbott:1982af,Preskill:1982cy,Dine:1982ah,Davis:1986xc}. 
From the experimental side, ALPs can be probed by cosmological and astrophysical  searches~\cite{ADMX:2003rdr,Barbieri:2016vwg,Caldwell:2016dcw,Zioutas:1998cc,Irastorza:2011gs,CAST:2017uph,Armengaud:2014gea,VanBibber:1987rq,Bahre:2013ywa,OSQAR:2015qdv,Arvanitaki:2014dfa}, beam-dump experiments~\cite{Alekhin:2015byh,Dobrich:2015jyk,Shan:2024pcc}, 
at colliders~\cite{Bauer:2017ris,Brivio:2017ije} and through a plethora of rare processes~\cite{Bauer:2021mvw,Gavela:2019wzg,Cornella:2019uxs}.
From the theoretical viewpoint, it is customary to describe the leading-order ALP interactions with SM particles via effective dimension-5 operators~\cite{Georgi:1986df}. Such an Effective Field Theory (EFT) approach allows to capture general features of broad classes of models without relying on specific ultraviolet completions. 
Physical observables are then obtained by computing matrix elements of the ALP EFT Lagrangian at those energy scales $E \ll \Lambda$ that are accessible by experiments. 
As a result, a crucial ingredient to make theoretical predictions is to run the ALP Lagrangian from the scale $\Lambda$ down to the experimental scale $E$.
This goal can be achieved by evaluating the anomalous dimension matrix of the ALP effective operators.

The renormalization group equations (RGEs) of the Standard Model EFT extended with a CP-odd ALP have been already computed at one-loop accuracy up to dimension-5 operators using diagrammatic methods~\cite{Bauer:2020jbp,Chala:2020wvs}. 
Instead, the case of ALPs with both CP even and odd components leads to CP-violating effects which have
been investigated in~\cite{Marciano:2016yhf,DiLuzio:2020oah}.
The corresponding RGEs of such a CP violating ALP framework
have been derived in~\cite{DasBakshi:2023lca}. 

Quite recently, anomalous dimensions have been evaluated  
through on-shell and unitarity-based techniques for scattering amplitudes~\cite{Caron-Huot:2016cwu,EliasMiro:2020tdv,Baratella:2020lzz,Jiang:2020mhe,Bern:2020ikv,Baratella:2020dvw,AccettulliHuber:2021uoa,EliasMiro:2021jgu,Baratella:2022nog,Machado:2022ozb}.
This approach, rooted in recent developments in the context of generalized unitarity \cite{Arkani-Hamed:2008owk, Huang:2012aq, Cheung:2015aba} and of the study of the dilatation operator in $\mathcal{N}=4$ super Yang-Mills theories \cite{Zwiebel:2011bx,Wilhelm:2014qua, Nandan:2014oga, Koster:2014fva, Brandhuber:2014pta, Brandhuber:2015boa, Loebbert:2015ova, Brandhuber:2016fni}, is particularly suited to unveil hidden structures with the emergence of zeros in the anomalous dimension matrix.
The origin of these vanishing elements is a direct consequence of selection rules~\cite{Elias-Miro:2014eia} based on operator lengths~\cite{Bern:2019wie}, helicity~\cite{Cheung:2015aba}, and angular momentum~\cite{Jiang:2020rwz}.\footnote{See also~\cite{Chala:2023jyx,Chala:2023xjy} for other zeros stemming from arguments based on the positivity of the $S$-matrix.}
%
Remarkably, anomalous dimensions can be 
related to the discontinuities of form factors of EFT operators,
therefore, they can be efficiently extracted from
unitarity cuts, evaluated via phase-space integrals ~\cite{Caron-Huot:2016cwu}.
This method has been proven to be
particularly promising for computing anomalous dimensions beyond one-loop order
\cite{Caron-Huot:2016cwu,Bern:2020ikv,EliasMiro:2021jgu}.

First studies concerned anomalous dimension matrices of non-renormalizable massless theories including mixing effects among operators of the same dimension~\cite{Caron-Huot:2016cwu}.
More recent studies have also systematically considered the inclusion of leading mass effects in this framework, which are extremely relevant in several EFT extensions of the SM~\cite{Bresciani:2023jsu}.
In particular, leading mass effects can be included in the massless limit~\cite{Bresciani:2023jsu} by exploiting the Higgs low-energy theorem~\cite{Ellis:1975ap,Shifman:1979eb,Shadmi:2018xan,Durieux:2019eor,Durieux:2020gip,Balkin:2021dko,Bertuzzo:2023slg}. 

The aim of this paper is to apply the above method~\cite{Caron-Huot:2016cwu,Bresciani:2023jsu} to the one-loop renormalization of CP violating 
ALP theories~\cite{Marciano:2016yhf,DiLuzio:2020oah} up to the phenomenologically most relevant dimension-6 operators, therefore reproducing and extending previous results~\cite{DasBakshi:2023lca}.
%
An extensive derivation of the relevant anomalous dimension matrix will be carried out at one-loop order both with standard techniques and through on-shell methods, aiming to show the strength of the latter approach, which drastically reduces the complexity of standard loop calculations. The relevant phase space cut-integrals will be evaluated by different parameterizations, both by angular integration~\cite{Zwiebel:2011bx,Caron-Huot:2016cwu,EliasMiro:2020tdv,Baratella:2020lzz,Bern:2020ikv,Baratella:2020dvw}, and using Stokes' theorem~\cite{Mastrolia:2009dr,Jiang:2020mhe,Shu:2021qlr,Baratella:2021guc}, 
for cross checks, as well as to highlight the strengths of the various approaches.

The paper is organized as follows. In Section~\ref{sec:method}, we summarize the 
on-shell method~\cite{Caron-Huot:2016cwu} and present different parameterizations to evaluate phase space integrals.
In Section~\ref{sec:ALPEFT}, we introduce the EFT for axion-like particles and in Section~\ref{sec:uvanomalousdim} we report a detailed derivation of the corresponding anomalous dimensions.
In Section~\ref{sec:comparison}, we compare our results as obtained with on-shell and standard methods.
Section~\ref{sec:conclusions} is dedicated to our conclusions.
In Appendix~\ref{app:notation} we report our notation and conventions and, finally, tree amplitudes and infrared anomalous dimensions of ALP operators are given in Appendix~\ref{app:Amplitudes} and \ref{app:IRanomalousdimensions}, respectively.

